% ---------------------------------------------------------------------------
% Chapter 5: Experimental Study
% MSc dissertation: Multi-Objective Evolutionary Design of Unlearning
% Strategies (MED-US)
%
% Evidence base: results/writeup_package/, results/literature_alignment/,
% results/analysis/class_structure/, results/search/plan_a_*/summary.json.
% Every numeric claim in this chapter is traced in
% dissertation/notes/chapter5_evidence_map.md.
%
% Packages assumed by this chapter: booktabs, graphicx, siunitx (optional),
% amsmath, csvsimple (optional, if tables are loaded directly from CSV).
% ---------------------------------------------------------------------------

\chapter{Experimental Study}
\label{ch:experimental-study}

\section{Introduction}
\label{sec:exp-intro}

This chapter reports the empirical behaviour of MED-US across ten independent
experiments, one for each CIFAR-10 class, together with the reference models they
are measured against, the benchmark they are compared with, and a separate
post-search refinement variant.

The chapter is organised so that measurement precedes interpretation.
Sections~\ref{sec:exp-setup} to \ref{sec:exp-reference} establish the setting and
the gold standard; Sections~\ref{sec:exp-pure} to \ref{sec:exp-operators} report
results against the research questions posed in Chapter~\ref{ch:introduction};
Section~\ref{sec:exp-truck} is a case study of the single class on which the
method performs worst; and Section~\ref{sec:exp-summary} consolidates the
findings.

Two commitments shape the presentation. First, the pure gradient-free method and
the hybrid refinement are reported separately throughout, because they make
different assumptions and only one is comparable with the gradient-free
literature. Second, results that do not favour the method are given the same
space as those that do: the forgetting performance reported below is substantially
weaker than the published state of the art, and the class on which the method
performs least well is given a dedicated section.

\section{Experimental Setup}
\label{sec:exp-setup}

\subsection{Hardware and software}

All searches, refinements, analyses and figures were produced on a single
workstation with an NVIDIA GeForce GTX 1650 (4.0\,GB VRAM, compute capability
7.5), an AMD Ryzen 5 3550H processor (4 cores, 8 threads) and 13.9\,GB of system
memory, running Windows 11. The software stack was Python 3.11, PyTorch 2.5.1
with CUDA 12.1, and torchvision 0.20.1.

Automatic mixed precision was deliberately disabled. Measured on this hardware
with ResNet-18 and CIFAR-10 at batch size 64, fp32 throughput was 440 samples per
second with a finite loss, while fp16 achieved 165 samples per second and produced
a NaN loss inside \texttt{layer4}. The GTX 1650 is the one Turing part without
tensor cores, so half precision has no fast path on it, and all experiments ran in
fp32.

The retain-only reference models were trained separately on Kaggle using NVIDIA
Tesla T4 accelerators, since ten full training runs were not practical on the
local GPU within the project timeline. That training is described in
Section~\ref{sec:exp-reference}.

\subsection{Determinism and reproducibility}

Deterministic algorithms are enforced inside the evaluator, so that evaluating
the same chromosome twice yields identical objective values. This property makes
the objective cache sound, since distinct genomes that decode to the same strategy
can be served from cache without altering the search, and it also permits a
selected candidate to be reconstructed from its stored chromosome and re-measured,
a check reported in Section~\ref{sec:exp-truck}.

Every search used seed 42. Search state, including the random number generator
stream, was checkpointed once per generation so that an interrupted run could be
resumed as the same experiment rather than as a new one.

\section{Dataset, Model, and Forgetting Protocol}
\label{sec:exp-data}

\subsection{Class-level unlearning on CIFAR-10}

The task in this dissertation is \emph{class-level} unlearning. Given a model
trained on all ten CIFAR-10 classes, the objective is to produce a model that
behaves as though one entire class had never been present in its training data,
while its behaviour on the remaining nine classes is preserved.

This differs from the instance-level setting, in which the forget set is a random
subset of training images drawn from across all classes. The distinction is not
cosmetic: as established in Chapter~\ref{ch:medus}, an instance-level forget set
shares its features, and therefore its weights, with the retained data, leaving no
forget-specific structure for a weight edit to remove.

\subsection{Notation and splits}

\begin{description}
  \item[$D_f$] The forget set: every image of the single target class. In
    CIFAR-10 this is 5{,}000 training images and 1{,}000 held-out test images.
  \item[$D_r$] The retain set: every image of the remaining nine classes. This is
    45{,}000 training images and 9{,}000 held-out test images.
  \item[$W_0$] The original model, trained on the complete CIFAR-10 training set
    and therefore on $D_f$ as well as $D_r$. There is one $W_0$, shared by all ten
    experiments.
  \item[$W_{ref}$] The retain-only reference model, trained from scratch on
    $D_{r,\text{train}}$ alone. It has never seen a single image of the target
    class. There are ten of these, one per class.
  \item[$C^{*}$] The candidate strategy selected from the Pareto front produced by
    the search, and the model that strategy produces when applied to $W_0$.
\end{description}

Because CIFAR-10 is balanced, these counts are identical for every class. The
split requires no random seed and no stratification, because no sampling is
involved: the class label determines the partition entirely. Splits are stored as
version-controlled index lists, and every split inside every imported reference
bundle was compared byte for byte against the local split before the corresponding
checkpoint was accepted.

\subsection{Model}

All experiments use ResNet-18 in its CIFAR variant, with a single
$3\times3$ stride-1 stem convolution and no initial max-pooling layer. The
torchvision stem, designed for $224\times224$ inputs, downsamples by a factor of
four before the first residual stage and discards most of the spatial information
in a $32\times32$ image. The CIFAR variant is the convention used by the line of
work from which this project's operators are drawn
[CITE: Fan et al., year]; [CITE: Kodge et al., year].

The network has 11{,}173{,}962 trainable parameters distributed across six layer
groups. The distribution is highly uneven: \texttt{layer4} alone holds
8{,}393{,}728 parameters, or 75.12\% of the model. This concentration is the
reason the edit-cost objective introduced in Chapter~\ref{ch:medus} is normalised
rather than absolute.

$W_0$ was trained for 200 epochs with SGD at learning rate 0.1, momentum 0.9 and
weight decay $5\times10^{-4}$, with the learning rate decayed by a factor of ten
every 40 epochs. It reaches 94.79\% test accuracy.

\section{Reference Model Validation}
\label{sec:exp-reference}

\subsection{Why the reference matters}

$W_{ref}$ is the empirical target of the whole method. The goal is not zero
accuracy on the forget class: a model driven to zero has learned to actively avoid
the correct answer, which is itself a detectable signature, whereas a model that
never saw the class misclassifies it as any naive model would. The target is
therefore the \emph{gap} to $W_{ref}$, and $f_1$ measures that gap as a divergence
between predictive distributions. Since an unsound reference would invalidate
every result that follows, the ten reference models were validated before any
search was run.

\subsection{Training protocol}

Each $W_{ref}$ was trained from scratch on its own $D_{r,\text{train}}$ of 45{,}000
images, for 200 epochs, at seed 42, using the same optimiser settings as $W_0$.
Training took approximately 2.32 hours per model on a Tesla T4.

Checkpoint selection used $D_{r,\text{test}}$ accuracy, with $D_{r,\text{test}}$
loss as tie-breaker. Forget-class test accuracy was logged every epoch as a
diagnostic but never influenced selection. This matters more than it may appear.
Selecting on full-test accuracy would be both diluted, since one tenth of the test
set belongs to a class the reference never trained on, and actively backwards: the
forget-class logit is never positively trained but still fluctuates, so an epoch
that happens to place a few more forget-class images into that class would score
higher. Such a rule would reward the reference for recognising precisely the thing
it must never have seen.

Reference training was distributed across three people and seven Kaggle output
bundles, with one class trained locally. Each imported bundle was validated by an
automated checker before installation, and each checkpoint's SHA-256 hash was
recorded.

\subsection{Validation outcome}

\begin{table}[htbp]
  \centering
  \caption{Retain-only reference models for all ten CIFAR-10 classes. Each model
  was trained on $D_{r,\text{train}}$ only. Forget-class test accuracy is 0.0000
  for all ten, which is the correctness condition rather than a result.}
  \label{tab:reference-validation}
  % PLACEHOLDER: reference model validation table
  % SOURCE: results/writeup_package/reference_model_validation_table.csv
  % COLUMNS: class_id, class_name, verdict, selected_epoch, D_f_test_acc,
  %          D_r_test_acc, full_test_acc, sha256_prefix, source_zip
  % Rendered version with commentary:
  %          results/writeup_package/reference_model_validation_table.md
\end{table}

All ten models passed validation. Every one reaches a forget-class test accuracy
of exactly 0.0000, which is the required correctness condition: a valid
retain-only reference should not recognise held-out target-class images as the
forgotten class.

Retain-set test accuracy across the ten references is
$0.9506 \pm 0.0057$, ranging from 0.9423 for horse to 0.9612 for cat. This spread
is narrow, which is important because it means no class has an unusually weak or
unusually strong gold standard that could distort its unlearning result. The
comparison across classes reported in Section~\ref{sec:exp-pure} is therefore a
property of the method rather than of the references.

Full-test accuracy sits near 0.85 for every reference, because one class in ten is
held out by construction. This is not a utility number and should not be read as
one.

\section{Pure MED-US Results}
\label{sec:exp-pure}

\subsection{What is being reported}

Every result in this section is produced by \emph{pure} MED-US: an evolutionary
search over layer-wise weight-editing strategies in which no gradient is ever
computed. Connection importance is estimated from forward passes alone, and no
model in this section has had a gradient step applied to it.

For each class, NSGA-II ran with a population of 10 for 50 generations, giving a
budget of 510 chromosome evaluations. The ten members of the final non-dominated
front were then re-measured at full fidelity on the complete evaluation sets. A
single selection rule was applied uniformly to all ten classes: $C^{*}$ is the
front member maximising the composite $\mathrm{ACC}_r \times (1 - \mathrm{ACC}_f)$,
which is the metric function used by the anchor study
[CITE: Kodge et al., year] rather than one devised here.

\subsection{Search behaviour}

All ten searches completed with zero failed evaluations and produced
non-dominated fronts of ten members each. Real evaluations per class ranged from
212 to 319, with 191 to 298 cache hits, giving cache hit rates between 37\% and
58\%. These hits follow directly from the latent genes in the encoding: inactive
layer groups retain their operator and intensity values, so distinct genomes
decode to the same strategy and can be served from cache.

\subsection{Per-class results}

\begin{table}[htbp]
  \centering
  \caption{Pure MED-US across all ten CIFAR-10 classes. Every row is gradient-free
  weight surgery. $C^{*}$ was selected by one rule applied uniformly: the front
  member maximising the composite $\mathrm{ACC}_r \times (1 - \mathrm{ACC}_f)$.}
  \label{tab:pure-ten-class}
  % PLACEHOLDER: pure MED-US 10-class table
  % SOURCE: results/writeup_package/pure_medus_10_class_table.csv
  % COLUMNS: class_id, class_name, front_position, operators, ACC_r, ACC_f,
  %          composite, MIA, selectivity_S, f1_js, f2_retain_train_loss,
  %          f3_edit_cost, search_minutes
  % Rendered version: results/writeup_package/pure_medus_10_class_table.md
\end{table}

Across the ten classes, pure MED-US achieves a retain accuracy of
$93.84 \pm 1.15$, a forget accuracy of $12.55 \pm 11.57$, a composite of
$82.09 \pm 11.00$ and an anchor membership-inference score of $92.54 \pm 8.62$.
All values are percentages, and all are means over the ten per-class experiments.

The $\pm$ figures require a caution that is repeated in
Section~\ref{sec:exp-threats-forward}. They are class-to-class standard
deviations over ten experiments at a single seed. They are \emph{not} run-to-run
variance, and they should not be read as uncertainty estimates.

\subsection{The shape of the result}

The most informative feature of Table~\ref{tab:pure-ten-class} is not either mean
but the difference in spread between them.

Retention is stable. Retain accuracy varies by 3.10 percentage points across all
ten classes, from 92.52 for frog to 95.62 for dog, with a standard deviation of
1.15. The search does not damage the nine retained classes unpredictably.

Forgetting is not stable. Forget accuracy varies by 42.10 percentage points, from
0.00 for airplane to 42.10 for truck, with a standard deviation of 11.57, almost
as large as the mean itself. The cost of the method is concentrated almost
entirely in forgetting rather than in retention.

Airplane is the strongest single result and the only class to reach a forget
accuracy of exactly 0.00, matching the retain-only reference on the forgetting
axis, with a composite of 92.86 and a membership-inference score of 100.00. Truck
is the weakest at 42.10, and is analysed separately in
Section~\ref{sec:exp-truck}. The gap between truck and the next worst class, bird
at 17.90, exceeds the gap between bird and the best non-airplane class.

\subsection{Selectivity}

Selectivity $S$, the ratio of forget loss gained to retain loss paid, ranges from
93.99 for truck to 4427.91 for dog, with a mean of $1053.88 \pm 1413.47$. The
distribution is strongly skewed, and the mean should be read alongside the range
rather than on its own.

The relevant comparison is with the instance-level setting described in
Chapter~\ref{ch:introduction}, where the highest selectivity observed across
10{,}534 evaluated strategies was 1.158. Every one of the ten class-level results
reported here exceeds that figure by at least an order of magnitude, and the
strongest exceeds it by more than three. Damage under class-level unlearning is
selective in a way that damage under instance-level unlearning was not.

\begin{figure}[htbp]
  \centering
  % PLACEHOLDER: class_structure_analysis.png
  % PATH: results/writeup_package/figures/class_structure_analysis.png
  \caption{Class structure explains the regime, not the ranking. Panel A: every
  CIFAR-10 class places between 84.1\% and 91.2\% of channels above the noise
  floor, against 0.55\% for an instance-level forget set where chance alone gives
  1.00\%. Panel B: that structure does not predict how hard a class is to forget.
  Panel C: cosine similarity between per-class channel-contrast vectors recovers
  the semantic grouping without being told it.}
  \label{fig:class-structure}
\end{figure}

Figure~\ref{fig:class-structure} accounts for this regime change and
simultaneously rules out the most obvious explanation for the per-class variation
within it. Panel A confirms the precondition: every class places between 84.1\%
and 91.2\% of channels above the noise floor of a null control, against 0.55\% for
an instance-level forget set, where the null control gives 1.00\% by construction.

Panel B tests whether the amount of structure predicts how difficult a class is to
forget. It does not. The Pearson correlation between median signal-to-noise ratio
and pure forget accuracy across the ten classes is $-0.04$. Truck is sixth of ten
on median SNR and fifth on channels above the floor, yet it is by far the worst on
forgetting, while automobile has the least structure of any class and forgets 3.3
times better than truck.

This is a negative result and it is reported as one. The measurement accounts for
the \emph{regime}, in that class-level forget sets are amenable to weight editing
where instance-level forget sets are not, but it does not account for the
\emph{ranking} within that regime.

\section{Benchmark Comparison}
\label{sec:exp-benchmark}

\subsection{The comparison and its limits}

The anchor study for this work is [CITE: Kodge et al., year], whose primary
setting is identical to this one: CIFAR-10, ResNet-18, single-class forgetting,
gradient-free, and deployable without retraining. Its Table~1 reports six
competing methods aggregated as a mean over all ten target classes, which is why
all ten classes were swept here rather than a single representative class.

One limitation must be stated before any number is compared. \textbf{The baseline
rows are as reported in the published literature. No published unlearning baseline
was re-implemented or re-run in this project.} The comparison therefore inherits
every difference between two codebases, including data augmentation,
membership-inference attack details, checkpoint selection and evaluation subsets.

\begin{table}[htbp]
  \centering
  \caption{Benchmark comparison, CIFAR-10 and ResNet-18, mean over all ten target
  classes. Rows marked as reported are transcribed from the anchor study and were
  not re-measured in this project. Rows marked as measured were produced by the
  harness described in this chapter.}
  \label{tab:benchmark}
  % PLACEHOLDER: benchmark comparison table
  % SOURCE: results/writeup_package/benchmark_comparison_table.csv
  % COLUMNS: method, source, measured_in_this_harness, gradient_free,
  %          ACC_r_mean, ACC_r_std, ACC_f_mean, ACC_f_std, MIA_mean, MIA_std
  % Rendered version: results/writeup_package/benchmark_comparison_table.md
\end{table}

\begin{figure}[htbp]
  \centering
  % PLACEHOLDER: benchmark_comparison.png
  % PATH: results/writeup_package/figures/benchmark_comparison.png
  \caption{Retention is competitive; forgetting is not. Each metric occupies its
  own panel on its own scale, so the panels are not comparable with one another.
  The anchor values are as reported and were not re-measured here. The dashed line
  marks retraining from scratch.}
  \label{fig:benchmark}
\end{figure}

\subsection{What the comparison shows}

\textbf{Retention is competitive.} Pure MED-US holds retain accuracy at 93.84,
which is 0.35 percentage points below the anchor's reported $94.19 \pm 0.50$ and
sits inside the range of the published field.

\textbf{Forgetting is not.} Pure MED-US reaches a forget accuracy of 12.55,
against the anchor's reported $0.03 \pm 0.09$. The gap is 12.52 percentage points.
This is the principal negative finding of the study and it is stated without
qualification: on the metric that defines the task, the method reported here
performs substantially worse than the published state of the art. It is also worse
than the nearest
published neighbour, UNSIR at $10.89 \pm 8.79$
[CITE: Tarun et al., year], although that method shares this one's characteristic
of wide per-class spread rather than uniform behaviour.

\textbf{Privacy is competitive on the anchor's own metric.} The
membership-inference score is 92.54, against the anchor's reported
$95.50 \pm 14.23$. Section~\ref{sec:exp-summary} returns to how much weight that
comparison can carry.

\subsection{Evidence that the two harnesses measure the same quantities}

The comparison is indirect, but it is not unsupported. Both harnesses measure two
models in common: the original model and the retraining gold standard. This
project's own measurements of those two models land within 0.10 and 0.25
percentage points of retain accuracy of the corresponding published rows, and
reproduce their forget accuracy and membership-inference values to within 0.10 and
0.03 respectively. That agreement on the shared baselines is the strongest
available evidence that the two implementations measure the same quantities. It
remains indirect, and does not substitute for re-running the baseline methods.

\section{Hybrid BN-Frozen Refinement Results}
\label{sec:exp-hybrid}

\subsection{A separate method, not a variant}

The results in this section were produced by a different method from the one in
Section~\ref{sec:exp-pure}, and are reported separately for that reason. The
hybrid takes the pure $C^{*}$ and applies two gradient steps outside the
evolutionary search: one clipped gradient-ascent step on $D_f$, followed by one
repair step of gradient descent on $D_r$. Batch normalisation is frozen throughout.
\textbf{The hybrid is therefore not gradient-free, and it is not a like-for-like
comparator for the gradient-free benchmark in Section~\ref{sec:exp-benchmark}.}

\subsection{Why batch normalisation must be frozen}

An earlier, unfrozen version of this refinement was initially recorded as a
success. It was not one. The model was left in training mode, and the eight
batches of $D_r$ used for the repair step silently re-estimated the batch
normalisation running statistics, undoing the operator edit, while every
weight-based guard, parameter movement, edit cost and retain accuracy alike,
continued to report success.

The failure was invisible to every check then in place, because batch
normalisation running statistics are buffers rather than parameters, so an edit
undone through them leaves no trace in any weight-based measurement. The remedy
was to hold the model in evaluation mode for both steps and to add an explicit
acceptance check on buffer movement.

\subsection{Acceptance criteria and outcome}

Each refinement had to pass six checks before its checkpoint was retained: forget
accuracy improved on $D_{f,\text{test}}$; retain test accuracy dropped by no more
than 0.010; no utility collapse, with retain losses no more than 1.25 times their
previous value; edit cost at most 0.3; parameter movement at most 0.0400; and
batch normalisation buffers unchanged.

\begin{table}[htbp]
  \centering
  \caption{Refinement acceptance record. Nine attempts, nine accepted, zero
  rejected. Airplane was not attempted because its pure forget accuracy is already
  0.00.}
  \label{tab:refinement-acceptance}
  % PLACEHOLDER: refinement acceptance table
  % SOURCE: results/writeup_package/refinement_acceptance_table.csv
  % COLUMNS: class_id, class_name, attempted, accepted, parameter_movement,
  %          movement_budget, buffer_movement, batchnorm_frozen,
  %          batchnorm_counters_changed, retain_test_drop, and six check_* columns
  % Rendered version: results/writeup_package/refinement_acceptance_table.md
\end{table}

Nine classes were eligible and all nine were accepted, with none rejected.
Airplane was deliberately not attempted: its pure forget accuracy is already 0.00
with a membership-inference score of 100.00, so a forgetting step has nothing to
improve, and its row in the hybrid table is the unchanged pure result.

Batch normalisation buffer movement was exactly 0.000000 on every accepted
refinement, with zero counter changes. Parameter movement ranged from 0.000303 for
cat to 0.000420 for frog against a budget of 0.0400, so every accepted refinement
used between 0.8\% and 1.1\% of the movement permitted.

\begin{table}[htbp]
  \centering
  \caption{Hybrid MED-US across all ten classes. Each refined row is the pure
  $C^{*}$ followed by one gradient-ascent step on $D_f$ and one repair step on
  $D_r$, with batch normalisation frozen. Airplane is an unmodified pure result.}
  \label{tab:hybrid-ten-class}
  % PLACEHOLDER: hybrid MED-US 10-class table
  % SOURCE: results/writeup_package/hybrid_medus_10_class_table.csv
  % COLUMNS: class_id, class_name, source, refinement_status, operators, ACC_r,
  %          ACC_f, composite, MIA, selectivity_S, parameter_movement,
  %          buffer_movement
  % Rendered version: results/writeup_package/hybrid_medus_10_class_table.md
\end{table}

Across the ten classes the hybrid achieves a retain accuracy of
$93.72 \pm 1.12$, a forget accuracy of $7.55 \pm 8.87$, a composite of
$86.66 \pm 8.52$ and a membership-inference score of $95.05 \pm 6.08$.

\section{Pure vs Hybrid Trade-off Analysis}
\label{sec:exp-tradeoff}

\subsection{The aggregate trade}

\begin{table}[htbp]
  \centering
  \caption{Pure against hybrid MED-US, ten-class means. These are two different
  methods and this table compares them; it does not merge them.}
  \label{tab:pure-vs-hybrid}
  % PLACEHOLDER: pure vs hybrid summary table
  % SOURCE: results/writeup_package/pure_vs_hybrid_summary_table.csv
  % COLUMNS: metric, n, pure_mean, pure_std, hybrid_mean, hybrid_std, delta_mean
  % The rendered version also carries the per-class delta table:
  %          results/writeup_package/pure_vs_hybrid_summary_table.md
\end{table}

The refinement reduces mean forget accuracy by 5.00 percentage points, from 12.55
to 7.55, at a cost of 0.12 percentage points of mean retain accuracy. The
composite rises by 4.58 and the membership-inference score by 2.51.

Nine of the ten classes improved on the composite. The tenth is the airplane
no-op, not a regression. No class deteriorated on any headline metric.

\begin{figure}[htbp]
  \centering
  % PLACEHOLDER: pure_vs_hybrid_acc_f_by_class.png
  % PATH: results/writeup_package/figures/pure_vs_hybrid_acc_f_by_class.png
  \caption{Forget accuracy falls for every refined class. Classes are ordered by
  pure forget accuracy, worst first, and the same order is used in
  Figures~\ref{fig:pvh-accr} and \ref{fig:pvh-composite} so that the three can be
  read across. Airplane is a deliberate no-op, so its two markers coincide. The
  hybrid applies two gradient steps outside the search and is not gradient-free.}
  \label{fig:pvh-accf}
\end{figure}

\begin{figure}[htbp]
  \centering
  % PLACEHOLDER: pure_vs_hybrid_acc_r_by_class.png
  % PATH: results/writeup_package/figures/pure_vs_hybrid_acc_r_by_class.png
  \caption{Retention is essentially unchanged by refinement. Note the axis: the
  entire ten-class range spans about three percentage points and no single class
  moves by more than 0.25. Markers rather than bars are used so that a non-zero
  axis is legitimate.}
  \label{fig:pvh-accr}
\end{figure}

\begin{figure}[htbp]
  \centering
  % PLACEHOLDER: pure_vs_hybrid_composite_by_class.png
  % PATH: results/writeup_package/figures/pure_vs_hybrid_composite_by_class.png
  \caption{Nine of ten classes improve on the composite. The tenth is the airplane
  no-op rather than a regression.}
  \label{fig:pvh-composite}
\end{figure}

\subsection{Who benefits}

The refinement helps the weakest classes most, which is why the standard deviation
narrows on every metric rather than only the mean shifting.

Truck gains 11.20 percentage points of forget accuracy, the largest absolute
improvement of any class, together with 10.24 on the composite and 9.50 on
membership inference. Ship is the runner-up at 9.30, moving from 14.00 to 4.70 and
finishing in a strong position with a composite of 90.05. At the other end, dog
gains only 1.40, because at 3.30 it had comparatively little left to gain.

The largest single-class retain cost is 0.2444 percentage points for cat. Frog is
the one class whose retain accuracy improves, by 0.0333. No class exceeded the
0.010 fractional retain-test drop enforced by the second acceptance check.

\subsection{Why the two are never merged}

The case for reporting the two together is not trivial, since the hybrid
outperforms the pure method on every headline metric. Three considerations make
the separation necessary.

First, they are different methods. Pure MED-US applies no gradient at any point;
the hybrid applies two.

Second, the anchor study's method is gradient-free. Only the pure results are a
like-for-like comparison with it, and placing a hybrid row in that comparison
would not be a fair one.

Third, reporting the hybrid's 7.55 as the result of MED-US would overstate what
the gradient-free method achieves by 5.00 percentage points of forget accuracy and
4.58 of composite.

The hybrid addresses a narrower and distinct question: what does one constrained
gradient step add to a gradient-free solution? The answer is 5.00 percentage
points of forget accuracy for 0.12 of retain accuracy. It also carries a
deployment cost, since it requires access to both $D_f$ and $D_r$ and an optimiser
step at unlearning time, which removes part of the practical argument for a
gradient-free method in the first place.

\section{Operator Frequency Analysis}
\label{sec:exp-operators}

\subsection{The observation}

\begin{figure}[htbp]
  \centering
  % PLACEHOLDER: operator_frequency_selected_cstar.png
  % PATH: results/writeup_package/figures/operator_frequency_selected_cstar.png
  \caption{MASK appears in the selected candidate for all ten classes, and no
  other operator appears in more than four. This is a search outcome rather than a
  design choice: all eight operators were available at equal cost in all ten
  independent runs.}
  \label{fig:operator-frequency}
\end{figure}

Across the ten selected candidates, the operator frequencies are MASK 10, CLIP 4,
DAMP 3, PRUNE 2, RANDOM\_PRUNE 2, RESET 2 and QUANTIZE 1, with NOISE appearing in
none. MASK is present in every one of the ten and is the only operator with that
property. Three classes, deer, dog and horse, select MASK alone, and five select
it either alone or with a single partner operator.

\subsection{Why this is a result rather than a setting}

All eight operators were available at equal cost in all ten searches, with
identical configuration and no per-class tuning, so nothing in the search biased
it towards MASK. The search converged on it independently ten times out of ten.

The convergence is consistent with the mechanism this project argues for. MASK is
the only editor-channel operator whose connection selector is the class-activation
contrast, scoring connections by
$|W_{ij}| \cdot (\mathrm{rms}^{f}_{j} - \mathrm{rms}^{r}_{j})$. It is therefore the
one operator in the library that can exploit the class-specific channels the
structure analysis in Figure~\ref{fig:class-structure} identified. PRUNE and
RANDOM\_PRUNE pin themselves to data-free selection rules, by magnitude and at
random respectively, and exist in the library precisely as controls: if a
forget-informed operator could not outperform them, the forget-informed component
would be doing nothing.

\subsection{The limit of the claim}

Ten runs at a single seed constitute convergent evidence, not proof. No operator
ablation was performed: the search was not re-run with MASK removed, nor with each
operator family in isolation, so it cannot be stated that MASK is necessary or
that the other operators are redundant. What can be stated is that, given a free
choice among eight operators in ten independent searches, the search selected MASK
every time.

\section{Truck Failure Analysis}
\label{sec:exp-truck}

\subsection{The failure}

Truck is the worst class in this study on every headline metric, both before and
after refinement.

Its pure forget accuracy is 42.10 against a ten-class mean of 12.55, more than
double the next worst class. Its composite is 53.80, against a next-worst of 78.10
for bird. Its membership-inference score is 69.60, against a next-worst of 91.10
for horse. Its selectivity is 93.99, the lowest of the ten and the only value below
100.

The refinement helps truck more than any other class, improving forget accuracy by
11.20 percentage points, and truck is still the worst class afterwards at 30.90
against a hybrid mean of 7.55. Refinement narrows the gap. It does not close it.

\subsection{Ruling out the obvious explanation}

The natural hypothesis is that truck carries less forget-specific structure than
other classes, leaving less for the operators to remove. Tested against the
class-structure measurement, this is false.

Truck sits mid-table on both structural measures, sixth of ten on median
signal-to-noise ratio and fifth of ten on the proportion of channels above the
noise floor, and the correlation across the ten classes between structure
magnitude and forget accuracy is a null $-0.04$. Automobile, which carries the
least forget-specific structure of any CIFAR-10 class, forgets 3.3 times better
than truck. Whatever makes truck difficult, it is not a shortage of structure.

\subsection{The destination of forgotten truck images}

To establish what does occur, all 1{,}000 held-out truck images were classified
under four models and their full predicted-class distributions recorded. This is
inference only; no model was trained, searched or refined for the analysis. The
pure $C^{*}$ has no stored checkpoint, since the search records genomes rather
than weights, so it was reconstructed by replaying its stored chromosome through
the same deterministic operators. The reconstruction reproduces the published
forget accuracy exactly, at 42.10 for the pure model and 30.90 for the hybrid,
confirming that the analysis describes the published models rather than near
neighbours of them.

\begin{figure}[htbp]
  \centering
  % PLACEHOLDER: truck_failure_analysis.png
  % PATH: results/writeup_package/figures/truck_failure_analysis.png
  \caption{Unlearning truck moves it toward automobile, and stalls. Predicted class
  of every held-out truck image under each of four models. A model that never saw a
  truck sends 68.40\% of them to automobile; pure MED-US sends only 16.70\% there
  and leaves 42.10\% still classified as truck.}
  \label{fig:truck-failure}
\end{figure}

The pattern is consistent across all four models. The original model classifies
95.40\% of truck images correctly as trucks. The retain-only reference, which has
never seen a truck, sends 68.40\% of them to automobile. Pure MED-US sends only
16.70\% to automobile and leaves 42.10\% still classified as truck. The hybrid
moves a further 11.20 percentage points out of the truck class, of which the
largest share arrives at automobile, raising that destination to 20.40\%.

The failure is therefore not a random scattering of the forgotten class across the
remaining nine. It is a \emph{partial movement along the truck to automobile
axis}: the reference model and the unlearned models share the same destination,
and the unlearned models simply do not travel as far.

\subsection{Corroboration from a second measurement}

An independent measurement supports the same reading. Cosine similarity between
the per-class channel-contrast vectors, shown in panel C of
Figure~\ref{fig:class-structure}, recovers the semantic grouping of CIFAR-10
without being given it: vehicles group with vehicles and animals with animals.
Within that structure, truck's nearest neighbour is automobile at 0.32, and the
relationship is mutual. Two measurements of very different kinds, one on the
activations of the original model and one on the predictions of four unlearned
models, identify the same pair.

\subsection{What this does not establish}

These observations do not establish that inter-class similarity causes the
difficulty, or that it predicts it. The correlation between each class's maximum
similarity to any other class and its forget accuracy is a null $-0.08$ across the
ten classes. More decisively, airplane is a direct counterexample: it has the
highest similarity to another class of any of the ten, at 0.41 with ship, and it is
the single class that reaches a forget accuracy of 0.00.

The defensible claim is therefore narrower. Truck's failure is stable and
reproducible rather than noise; it is not accounted for by the amount of
forget-specific structure the class carries; and it coincides with truck sharing
more of its structure with a retained class than with anything else. Why truck is
difficult remains an open question.

The case is reported in detail rather than in passing because a method with a
stable, identifiable and mechanistically characterised failure mode is of greater
practical value than one that fails unpredictably.

\section{Summary of Experimental Findings}
\label{sec:exp-summary}

\subsection{What was established}

\begin{enumerate}
  \item All ten retain-only reference models are valid, reaching forget-class test
    accuracy of 0.0000 with a retain-set test accuracy of $0.9506 \pm 0.0057$.
  \item Class-level forget sets carry forget-specific structure, with 84.1\% to
    91.2\% of channels above a null control's noise floor, against 0.55\% for the
    instance-level setting where chance gives 1.00\%.
  \item Pure gradient-free MED-US achieves $\mathrm{ACC}_r = 93.84 \pm 1.15$,
    $\mathrm{ACC}_f = 12.55 \pm 11.57$, composite $= 82.09 \pm 11.00$ and
    $\mathrm{MIA} = 92.54 \pm 8.62$ across ten classes, with zero failed searches.
  \item Retention is competitive with the published gradient-free state of the
    art, at 0.35 percentage points behind. Forgetting is not, at 12.52 percentage
    points worse.
  \item The hybrid refinement, which is not gradient-free, improves forget
    accuracy to $7.55 \pm 8.87$ for 0.12 percentage points of retain accuracy,
    with all nine attempts accepted and zero buffer movement throughout.
  \item MASK appears in all ten selected candidates and no other operator appears
    in more than four, without any operator ablation to establish necessity.
  \item Per-class difficulty is not predicted by the amount of forget-specific
    structure a class carries ($-0.04$), nor by its maximum similarity to another
    class ($-0.08$).
  \item Truck's failure is a partial movement along the truck to automobile axis,
    corroborated by two independent measurements, and its cause remains open.
\end{enumerate}

\subsection{Threats to validity}
\label{sec:exp-threats-forward}

Four limitations bear directly on how these results should be read.

\textbf{A single seed.} All results come from one seed, 42, for the original
model, all ten references and every search. Every $\pm$ in this chapter is
class-to-class spread across ten experiments, not run-to-run variance. Neither
search variance nor training variance was measured.

\textbf{No baseline was re-implemented.} The comparison in
Section~\ref{sec:exp-benchmark} is against published values. It does not support
any claim that MED-US outperforms or underperforms a specific method under
controlled conditions.

\textbf{Selection and reporting share a test set.} $C^{*}$ is chosen as the front
member maximising the composite computed on the same held-out sets on which the
result is subsequently reported. It is a best-of-ten selection on the reported
metric, and no separate selection split was used. The direction of the resulting
bias is optimistic, and its magnitude was not measured.

\textbf{The membership-inference metric may be saturated.} In the anchor's own
table, retraining scores exactly 100.00 and one competing method reaches a forget
accuracy of 0.00 with a membership-inference score of 0.00. A metric on which the
gold standard sits at the ceiling has limited discriminative power. This project's
own AUC-based attack, computed on the same models, sits far closer to chance than
an anchor score in the nineties would suggest. The two cannot be describing the
same underlying quantity, and the anchor score should not be read as strong
evidence of privacy.

\subsection{Cost}

The complete ten-class sweep required approximately 1.5 hours of search time in
total, ranging from 6.4 to 15.5 minutes per class, with a further 9.3 to 10.9
minutes per class for full-fidelity re-measurement of the front. End to end, a
single class costs roughly 30 to 35 minutes on a consumer GPU. By comparison,
training a single retain-only reference model took approximately 2.32 hours on a
datacentre accelerator. The method requires no retraining, no optimiser state and
no pass of gradient descent over the retain set.

These figures describe the shape of the cost rather than a competitive runtime
claim. The implementation is unoptimised: dataloader workers are disabled during
search, evaluation is capped at three batches, and the population is evaluated
serially rather than in parallel. A wall-clock comparison against published
runtimes would not be like-for-like.
